\documentclass[10pt]{article}
\usepackage[margin=0.8in]{geometry}
\usepackage[T1]{fontenc}
\usepackage{amsmath,amssymb,amsthm}
\usepackage{multicol}
\usepackage{hyperref}
\newtheorem{theorem}{Theorem}
\newtheorem{lemma}{Lemma}
\title{Cube-Sum Permutations: An 82-Identity Elementary Certificate}
\author{Qichao Wang}
\date{October 2026}

\begin{document}
\maketitle

\begin{abstract}
Let $P(n)$ count permutations $p$ of $\{1,\ldots,n\}$
for which every $i+p(i)$ is a perfect cube.
We prove $P(n)>0$ for every $n\ge 173$ from exactly
$82$ explicit endpoint-multiset identities, comprising
$1{,}030$ reflected integer intervals, plus an elementary
suffix-reflection rule and a cubic-gap inequality.
An explicit predecessor table explains all remaining
finite lengths without a matching search, a large
permutation archive, or dependent augmenting paths.
The finite case distinctions are \emph{not} eliminated:
this is a solver-free, hand-checkable certificate rather
than a case-free parameterized construction.
\end{abstract}

\section{Two elementary lemmas}

\begin{lemma}[Boundary cancellation]\label{lem:balance}
Let $0=b_0<b_1<\cdots<b_t=n$ and let
$c_1,\ldots,c_t$ be positive integer cubes.
Set $L_j=c_j-b_j$ and $R_j=c_j-b_{j-1}$.
Assume $1\le L_j<R_j\le n+1$.
Then $p(i)=c_j-i$ for $b_{j-1}<i\le b_j$
is a cube-sum permutation if and only if
\[
\boxed{\{1,R_1,\ldots,R_t\}_{\mathrm{multi}}
 =\{n+1,L_1,\ldots,L_t\}_{\mathrm{multi}}.}
\]
\end{lemma}
\begin{proof}
Each input block maps bijectively onto the integer
interval $[L_j,R_j)$. At an integer $x$, the jump
in the total interval-coverage count is the
multiplicity of $x$ as a left endpoint minus
its multiplicity as a right endpoint.
The displayed identity makes that jump $1$
at $x=1$, $-1$ at $x=n+1$, and zero everywhere
else. Thus coverage is exactly one at
every $x\in[1,n]$. Reversing this argument
proves necessity.
\end{proof}

\begin{lemma}[Growing cubic suffix]\label{lem:suffix}
Suppose $0\le m<n$, $m+n+1=c$ is a cube,
and a cube-sum permutation on $[1,m]$ exists
(the $m=0$ permutation is empty).
Keep the existing permutation on $[1,m]$
and set $p(i)=c-i$ for $m<i\le n$.
Then $p$ is a cube-sum permutation on $[1,n]$.
\end{lemma}
\begin{proof}
The new reflection is a bijection from
$[m+1,n]$ to itself, and each new sum is $c$.
\end{proof}

\section{An explicit 82-identity finite certificate}

The appendix records $82$ rows of the form
$n\; b_1:k_1\;\cdots\;b_t:k_t$.
The corresponding cubical sums are $c_j=k_j^3$
and the first boundary is $b_0=0$.
For every row, the boundaries increase,
$1\le L_j<R_j\le n+1$,
and the multiset equality in
Lemma~\ref{lem:balance} holds exactly.
These are $1{,}030$ simple integer-interval
identities. Exactly $22$ listed rows have
$n<173$ and exactly $60$ have $173\le n\le364$.

The following two lists give every omitted
base length without further matching data.
For any tabulated predecessor $m$ under
cube $c$, define $n=c-m-1$ and apply
Lemma~\ref{lem:suffix}. The lists are
ordered by cube, so every predecessor is
either a seed or was obtained in an earlier row.

\paragraph{First: 26 reflected lengths below 173.}
\begin{quote}
\footnotesize
\begin{verbatim}
cube     already constructed predecessors m
  8      0
 27      0,7
 64      0,7,19,26
125      0,7,19,26,37,44,56
216      44,56,63,66,68,80,82,85,87,
         100,103,105
\end{verbatim}
\end{quote}
Together with the $22$ lower seeds,
these supply $48$ distinct lengths below 173.

\paragraph{Second: 49 reflected lengths in
$[173,255]\cup[339,364]$.}
\begin{quote}
\footnotesize
\begin{verbatim}
cube     previously constructed predecessors m
216      0,7,19,26,37
343      0,87,98,100,103,105,110,112,
         115,116,117,118,119,120,121,122,
         123,124,125,126,128,129,130,131,
         132,133,135,147,149,150,151,152,
         155,156,159
512      147,149,150,151,152,155,156,
         159,171
\end{verbatim}
\end{quote}
The $60$ upper seed rows and these $49$
reflections are exactly the $109$ integers
in $[173,255]\cup[339,364]$.
There is no circular dependency:
each predecessor $m$ is strictly smaller
than its constructed length $n$.

For example, $211=343-131-1$,
$340=512-171-1$, and $342=343-0-1$.
The first example replaces a previously
necessary-looking $53$-segment independent
row by one suffix reflection.

\begin{lemma}[Complete finite interval]\label{lem:base}
For every integer $173\le n\le364$,
a cube-sum permutation of $[1,n]$ exists.
\end{lemma}
\begin{proof}
The explicit arithmetic rows and the two
predecessor tables establish the claim for
$[173,255]\cup[339,364]$. For each of the
$83$ remaining lengths $256\le n\le338$,
take $m=511-n\in[173,255]$.
Then $m<n$ and $m+n+1=512=8^3$,
so Lemma~\ref{lem:suffix} applies.
\end{proof}

\section{A purely symbolic infinite tail}

\begin{theorem}
For every integer $n\ge173$,
there is a permutation $p$ of $[1,n]$
such that $i+p(i)$ is a perfect cube
for all indices $i$.
\end{theorem}
\begin{proof}
The base interval is Lemma~\ref{lem:base}.
For $n\ge365$, assume the claim for
every smaller integer at least 173,
and let $k^3$ be the least integer cube
satisfying $k^3\ge n+174$.
Thus $m=k^3-n-1\ge173$.

If $365\le n\le555$, then $k=9$
and $k^3=729\le2n$. For $n\ge556$,
minimality of $k$ implies $k\ge10$
and $n\ge(k-1)^3-173$. Therefore
\[
\begin{aligned}
2n-k^3
&\ge 2(k-1)^3-346-k^3\\
&=k^3-6k^2+6k-348\\
&\ge112>0.
\end{aligned}
\]
The final polynomial equals $112$ at
$k=10$ and increases thereafter:
its first difference is $3k^2-9k+1>0$.
In either case $k^3\le2n$ and hence
$m<n$. Apply the induction hypothesis
at $m$ and Lemma~\ref{lem:suffix}.
\end{proof}

\paragraph{Logical scope.}
The theorem uses explicit finite
arithmetic equalities and short symbolic
lemmas. The optional script
\texttt{verify\_82\_seed\_proof.py}
checks the $82$ rows and finite closure
without matching or other external data.
It is not a premise for the proof.
A uniform, case-free formula replacing
the $82$ rows is \emph{not} claimed.

\clearpage
\appendix
\section{All 82 explicit integer-multiset witnesses}
\label{app:table}
Every token $b:k$ means $p(i)=k^3-i$
up to endpoint $b$. For each row, verify
the same endpoint balance of
Lemma~\ref{lem:balance}.

\begin{multicols}{2}
\scriptsize\ttfamily
\setlength{\parindent}{0pt}
\setlength{\parskip}{2pt}
% Exact 82 arithmetic identity rows, grouped into 22 lower and 60 middle seeds.
% Each n b:k token means p(i)=k^3-i up to the specified endpoint.
% Companion proof: COMPRESSED_FINITE_PROOF.md; source: finite_82_seed_atlas.txt.
\noindent \textbf{66}: 2:4\allowbreak\ 5:2\allowbreak\ 58:4\allowbreak\ 61:5\allowbreak\ 63:4\allowbreak\ 66:5\par
\noindent \textbf{82}: 2:4\allowbreak\ 5:2\allowbreak\ 21:3\allowbreak\ 42:4\allowbreak\ 61:5\allowbreak\ 63:4\allowbreak\ 82:5\par
\noindent \textbf{85}: 2:4\allowbreak\ 24:3\allowbreak\ 39:4\allowbreak\ 61:5\allowbreak\ 63:4\allowbreak\ 85:5\par
\noindent \textbf{100}: 2:4\allowbreak\ 24:3\allowbreak\ 61:5\allowbreak\ 63:4\allowbreak\ 100:5\par
\noindent \textbf{103}: 2:4\allowbreak\ 5:2\allowbreak\ 21:3\allowbreak\ 61:5\allowbreak\ 63:4\allowbreak\ 103:5\par
\noindent \textbf{116}: 2:4\allowbreak\ 3:3\allowbreak\ 4:2\allowbreak\ 8:3\allowbreak\ 9:5\allowbreak\ 10:3\allowbreak\ 11:5\allowbreak\ 15:3\allowbreak\ 16:5\allowbreak\ 17:3\allowbreak\ 18:5\allowbreak\ 22:3\allowbreak\ 23:5\allowbreak\ 24:3\allowbreak\ 61:5\allowbreak\ 63:4\allowbreak\ 100:5\allowbreak\ 101:6\allowbreak\ 102:5\allowbreak\ 106:6\allowbreak\ 107:5\allowbreak\ 108:6\allowbreak\ 109:5\allowbreak\ 113:6\allowbreak\ 114:5\allowbreak\ 115:6\allowbreak\ 116:5\par
\noindent \textbf{118}: 1:3\allowbreak\ 2:2\allowbreak\ 3:3\allowbreak\ 4:2\allowbreak\ 5:3\allowbreak\ 6:2\allowbreak\ 7:5\allowbreak\ 8:3\allowbreak\ 9:5\allowbreak\ 10:3\allowbreak\ 11:5\allowbreak\ 12:3\allowbreak\ 14:5\allowbreak\ 15:3\allowbreak\ 16:5\allowbreak\ 17:3\allowbreak\ 18:5\allowbreak\ 19:3\allowbreak\ 21:5\allowbreak\ 22:3\allowbreak\ 23:5\allowbreak\ 24:3\allowbreak\ 25:5\allowbreak\ 26:3\allowbreak\ 98:5\allowbreak\ 99:6\allowbreak\ 100:5\allowbreak\ 101:6\allowbreak\ 102:5\allowbreak\ 103:6\allowbreak\ 105:5\allowbreak\ 106:6\allowbreak\ 107:5\allowbreak\ 108:6\allowbreak\ 109:5\allowbreak\ 110:6\allowbreak\ 112:5\allowbreak\ 113:6\allowbreak\ 114:5\allowbreak\ 115:6\allowbreak\ 116:5\allowbreak\ 117:6\allowbreak\ 118:5\par
\noindent \textbf{119}: 5:3\allowbreak\ 7:5\allowbreak\ 12:3\allowbreak\ 14:5\allowbreak\ 19:3\allowbreak\ 21:5\allowbreak\ 26:3\allowbreak\ 98:5\allowbreak\ 103:6\allowbreak\ 105:5\allowbreak\ 110:6\allowbreak\ 112:5\allowbreak\ 117:6\allowbreak\ 119:5\par
\noindent \textbf{120}: 2:4\allowbreak\ 3:3\allowbreak\ 4:2\allowbreak\ 9:5\allowbreak\ 10:3\allowbreak\ 16:5\allowbreak\ 17:3\allowbreak\ 23:5\allowbreak\ 24:3\allowbreak\ 61:5\allowbreak\ 63:4\allowbreak\ 100:5\allowbreak\ 101:6\allowbreak\ 107:5\allowbreak\ 108:6\allowbreak\ 114:5\allowbreak\ 115:6\allowbreak\ 120:5\par
\noindent \textbf{121}: 2:4\allowbreak\ 3:3\allowbreak\ 9:5\allowbreak\ 10:3\allowbreak\ 16:5\allowbreak\ 17:3\allowbreak\ 23:5\allowbreak\ 24:3\allowbreak\ 61:5\allowbreak\ 63:4\allowbreak\ 100:5\allowbreak\ 101:6\allowbreak\ 107:5\allowbreak\ 108:6\allowbreak\ 114:5\allowbreak\ 115:6\allowbreak\ 121:5\par
\noindent \textbf{122}: 2:4\allowbreak\ 61:5\allowbreak\ 63:4\allowbreak\ 122:5\par
\noindent \textbf{123}: 2:4\allowbreak\ 31:5\allowbreak\ 32:4\allowbreak\ 61:5\allowbreak\ 63:4\allowbreak\ 92:5\allowbreak\ 93:6\allowbreak\ 122:5\allowbreak\ 123:6\par
\noindent \textbf{125}: 2:4\allowbreak\ 3:3\allowbreak\ 4:2\allowbreak\ 9:5\allowbreak\ 10:3\allowbreak\ 16:5\allowbreak\ 17:3\allowbreak\ 23:5\allowbreak\ 24:3\allowbreak\ 29:5\allowbreak\ 34:4\allowbreak\ 61:5\allowbreak\ 63:4\allowbreak\ 90:5\allowbreak\ 95:6\allowbreak\ 100:5\allowbreak\ 101:6\allowbreak\ 107:5\allowbreak\ 108:6\allowbreak\ 114:5\allowbreak\ 115:6\allowbreak\ 120:5\allowbreak\ 125:6\par
\noindent \textbf{126}: 2:4\allowbreak\ 5:2\allowbreak\ 28:5\allowbreak\ 35:4\allowbreak\ 61:5\allowbreak\ 63:4\allowbreak\ 89:5\allowbreak\ 96:6\allowbreak\ 119:5\allowbreak\ 126:6\par
\noindent \textbf{127}: 1:5\allowbreak\ 4:2\allowbreak\ 5:3\allowbreak\ 6:2\allowbreak\ 11:5\allowbreak\ 12:3\allowbreak\ 18:5\allowbreak\ 19:3\allowbreak\ 25:5\allowbreak\ 26:3\allowbreak\ 27:5\allowbreak\ 30:4\allowbreak\ 31:5\allowbreak\ 36:4\allowbreak\ 88:5\allowbreak\ 91:6\allowbreak\ 92:5\allowbreak\ 97:6\allowbreak\ 102:5\allowbreak\ 103:6\allowbreak\ 109:5\allowbreak\ 110:6\allowbreak\ 116:5\allowbreak\ 117:6\allowbreak\ 118:5\allowbreak\ 121:6\allowbreak\ 122:5\allowbreak\ 127:6\par
\noindent \textbf{129}: 3:5\allowbreak\ 4:2\allowbreak\ 8:3\allowbreak\ 11:5\allowbreak\ 15:3\allowbreak\ 18:5\allowbreak\ 22:3\allowbreak\ 25:5\allowbreak\ 30:4\allowbreak\ 33:5\allowbreak\ 38:4\allowbreak\ 86:5\allowbreak\ 91:6\allowbreak\ 94:5\allowbreak\ 99:6\allowbreak\ 102:5\allowbreak\ 106:6\allowbreak\ 109:5\allowbreak\ 113:6\allowbreak\ 116:5\allowbreak\ 121:6\allowbreak\ 124:5\allowbreak\ 129:6\par
\noindent \textbf{131}: 3:5\allowbreak\ 23:3\allowbreak\ 30:4\allowbreak\ 33:5\allowbreak\ 40:4\allowbreak\ 84:5\allowbreak\ 91:6\allowbreak\ 94:5\allowbreak\ 121:6\allowbreak\ 124:5\allowbreak\ 131:6\par
\noindent \textbf{132}: 3:5\allowbreak\ 4:2\allowbreak\ 22:3\allowbreak\ 30:4\allowbreak\ 33:5\allowbreak\ 41:4\allowbreak\ 83:5\allowbreak\ 91:6\allowbreak\ 94:5\allowbreak\ 121:6\allowbreak\ 124:5\allowbreak\ 132:6\par
\noindent \textbf{150}: 3:5\allowbreak\ 4:2\allowbreak\ 30:4\allowbreak\ 33:5\allowbreak\ 59:4\allowbreak\ 65:5\allowbreak\ 91:6\allowbreak\ 94:5\allowbreak\ 121:6\allowbreak\ 124:5\allowbreak\ 150:6\par
\noindent \textbf{151}: 3:5\allowbreak\ 30:4\allowbreak\ 33:5\allowbreak\ 60:4\allowbreak\ 64:5\allowbreak\ 91:6\allowbreak\ 94:5\allowbreak\ 121:6\allowbreak\ 124:5\allowbreak\ 151:6\par
\noindent \textbf{155}: 3:5\allowbreak\ 30:4\allowbreak\ 33:5\allowbreak\ 60:4\allowbreak\ 91:6\allowbreak\ 94:5\allowbreak\ 121:6\allowbreak\ 124:5\allowbreak\ 155:6\par
\noindent \textbf{156}: 3:5\allowbreak\ 4:2\allowbreak\ 30:4\allowbreak\ 33:5\allowbreak\ 59:4\allowbreak\ 91:6\allowbreak\ 94:5\allowbreak\ 121:6\allowbreak\ 124:5\allowbreak\ 156:6\par
\noindent \textbf{173}: 7:2\allowbreak\ 12:3\allowbreak\ 14:5\allowbreak\ 19:3\allowbreak\ 21:5\allowbreak\ 42:4\allowbreak\ 103:6\allowbreak\ 105:5\allowbreak\ 110:6\allowbreak\ 112:5\allowbreak\ 173:6\par
\noindent \textbf{174}: 3:5\allowbreak\ 4:2\allowbreak\ 22:3\allowbreak\ 30:4\allowbreak\ 33:5\allowbreak\ 41:4\allowbreak\ 91:6\allowbreak\ 94:5\allowbreak\ 121:6\allowbreak\ 124:5\allowbreak\ 174:6\par
\noindent \textbf{175}: 3:5\allowbreak\ 23:3\allowbreak\ 30:4\allowbreak\ 33:5\allowbreak\ 40:4\allowbreak\ 91:6\allowbreak\ 94:5\allowbreak\ 121:6\allowbreak\ 124:5\allowbreak\ 175:6\par
\noindent \textbf{176}: 7:2\allowbreak\ 12:3\allowbreak\ 14:5\allowbreak\ 17:4\allowbreak\ 19:3\allowbreak\ 24:5\allowbreak\ 42:4\allowbreak\ 103:6\allowbreak\ 105:5\allowbreak\ 110:6\allowbreak\ 115:5\allowbreak\ 166:6\allowbreak\ 169:7\allowbreak\ 176:6\par
\noindent \textbf{177}: 3:5\allowbreak\ 4:2\allowbreak\ 8:3\allowbreak\ 11:5\allowbreak\ 15:3\allowbreak\ 18:5\allowbreak\ 22:3\allowbreak\ 25:5\allowbreak\ 30:4\allowbreak\ 33:5\allowbreak\ 38:4\allowbreak\ 91:6\allowbreak\ 94:5\allowbreak\ 99:6\allowbreak\ 102:5\allowbreak\ 106:6\allowbreak\ 109:5\allowbreak\ 113:6\allowbreak\ 116:5\allowbreak\ 121:6\allowbreak\ 124:5\allowbreak\ 177:6\par
\noindent \textbf{179}: 3:5\allowbreak\ 11:3\allowbreak\ 15:4\allowbreak\ 23:3\allowbreak\ 30:4\allowbreak\ 33:5\allowbreak\ 40:4\allowbreak\ 48:6\allowbreak\ 52:4\allowbreak\ 91:6\allowbreak\ 94:5\allowbreak\ 121:6\allowbreak\ 124:5\allowbreak\ 163:6\allowbreak\ 167:7\allowbreak\ 175:6\allowbreak\ 179:7\par
\noindent \textbf{180}: 5:3\allowbreak\ 7:5\allowbreak\ 12:3\allowbreak\ 14:5\allowbreak\ 19:3\allowbreak\ 21:5\allowbreak\ 26:3\allowbreak\ 28:5\allowbreak\ 35:4\allowbreak\ 96:6\allowbreak\ 98:5\allowbreak\ 103:6\allowbreak\ 105:5\allowbreak\ 110:6\allowbreak\ 112:5\allowbreak\ 117:6\allowbreak\ 119:5\allowbreak\ 180:6\par
\noindent \textbf{181}: 3:5\allowbreak\ 4:2\allowbreak\ 29:5\allowbreak\ 30:4\allowbreak\ 33:5\allowbreak\ 34:4\allowbreak\ 91:6\allowbreak\ 94:5\allowbreak\ 95:6\allowbreak\ 120:5\allowbreak\ 121:6\allowbreak\ 124:5\allowbreak\ 181:6\par
\noindent \textbf{182}: 33:5\allowbreak\ 91:6\allowbreak\ 124:5\allowbreak\ 182:6\par
\noindent \textbf{184}: 7:2\allowbreak\ 31:4\allowbreak\ 32:6\allowbreak\ 56:4\allowbreak\ 159:6\allowbreak\ 183:7\allowbreak\ 184:6\par
\noindent \textbf{185}: 21:5\allowbreak\ 26:3\allowbreak\ 30:4\allowbreak\ 51:6\allowbreak\ 58:4\allowbreak\ 91:6\allowbreak\ 112:5\allowbreak\ 178:6\allowbreak\ 185:7\par
\noindent \textbf{188}: 2:5\allowbreak\ 5:4\allowbreak\ 7:2\allowbreak\ 30:4\allowbreak\ 32:6\allowbreak\ 61:4\allowbreak\ 91:6\allowbreak\ 93:5\allowbreak\ 154:6\allowbreak\ 157:7\allowbreak\ 159:6\allowbreak\ 188:7\par
\noindent \textbf{194}: 7:2\allowbreak\ 12:3\allowbreak\ 14:5\allowbreak\ 19:3\allowbreak\ 21:5\allowbreak\ 103:6\allowbreak\ 105:5\allowbreak\ 110:6\allowbreak\ 112:5\allowbreak\ 194:6\par
\noindent \textbf{197}: 8:5\allowbreak\ 18:3\allowbreak\ 55:6\allowbreak\ 63:4\allowbreak\ 91:6\allowbreak\ 99:5\allowbreak\ 182:6\allowbreak\ 190:7\allowbreak\ 197:6\par
\noindent \textbf{198}: 2:4\allowbreak\ 5:2\allowbreak\ 28:5\allowbreak\ 35:4\allowbreak\ 53:6\allowbreak\ 61:5\allowbreak\ 63:4\allowbreak\ 71:5\allowbreak\ 96:6\allowbreak\ 119:5\allowbreak\ 144:6\allowbreak\ 162:7\allowbreak\ 180:6\allowbreak\ 198:7\par
\noindent \textbf{199}: 2:4\allowbreak\ 5:2\allowbreak\ 16:4\allowbreak\ 22:6\allowbreak\ 41:4\allowbreak\ 47:6\allowbreak\ 58:4\allowbreak\ 61:5\allowbreak\ 63:4\allowbreak\ 66:5\allowbreak\ 149:6\allowbreak\ 168:7\allowbreak\ 174:6\allowbreak\ 193:7\allowbreak\ 199:6\par
\noindent \textbf{200}: 26:3\allowbreak\ 37:4\allowbreak\ 51:6\allowbreak\ 73:5\allowbreak\ 142:6\allowbreak\ 164:7\allowbreak\ 178:6\allowbreak\ 200:7\par
\noindent \textbf{201}: 1:3\allowbreak\ 12:4\allowbreak\ 25:3\allowbreak\ 37:4\allowbreak\ 50:6\allowbreak\ 62:5\allowbreak\ 63:4\allowbreak\ 141:6\allowbreak\ 164:7\allowbreak\ 177:6\allowbreak\ 190:7\allowbreak\ 201:6\par
\noindent \textbf{202}: 13:4\allowbreak\ 25:6\allowbreak\ 38:4\allowbreak\ 50:6\allowbreak\ 63:4\allowbreak\ 152:6\allowbreak\ 165:7\allowbreak\ 177:6\allowbreak\ 190:7\allowbreak\ 202:6\par
\noindent \textbf{203}: 12:4\allowbreak\ 25:6\allowbreak\ 26:3\allowbreak\ 37:4\allowbreak\ 51:6\allowbreak\ 62:4\allowbreak\ 152:6\allowbreak\ 164:7\allowbreak\ 178:6\allowbreak\ 189:7\allowbreak\ 203:6\par
\noindent \textbf{204}: 11:4\allowbreak\ 15:6\allowbreak\ 26:3\allowbreak\ 36:4\allowbreak\ 51:6\allowbreak\ 61:5\allowbreak\ 142:6\allowbreak\ 163:7\allowbreak\ 178:6\allowbreak\ 188:7\allowbreak\ 204:6\par
\noindent \textbf{205}: 2:4\allowbreak\ 5:2\allowbreak\ 21:3\allowbreak\ 42:4\allowbreak\ 46:6\allowbreak\ 61:5\allowbreak\ 63:4\allowbreak\ 78:5\allowbreak\ 137:6\allowbreak\ 169:7\allowbreak\ 173:6\allowbreak\ 205:7\par
\noindent \textbf{206}: 7:2\allowbreak\ 18:3\allowbreak\ 19:4\allowbreak\ 20:6\allowbreak\ 44:4\allowbreak\ 45:6\allowbreak\ 55:5\allowbreak\ 56:4\allowbreak\ 79:5\allowbreak\ 136:6\allowbreak\ 146:7\allowbreak\ 147:6\allowbreak\ 171:7\allowbreak\ 172:6\allowbreak\ 206:7\par
\noindent \textbf{228}: 2:4\allowbreak\ 3:2\allowbreak\ 4:3\allowbreak\ 5:2\allowbreak\ 22:6\allowbreak\ 23:3\allowbreak\ 41:5\allowbreak\ 58:4\allowbreak\ 61:5\allowbreak\ 63:4\allowbreak\ 101:5\allowbreak\ 114:6\allowbreak\ 132:7\allowbreak\ 149:6\allowbreak\ 228:7\par
\noindent \textbf{229}: 2:4\allowbreak\ 5:2\allowbreak\ 22:6\allowbreak\ 41:5\allowbreak\ 58:4\allowbreak\ 61:5\allowbreak\ 63:4\allowbreak\ 102:5\allowbreak\ 113:6\allowbreak\ 132:7\allowbreak\ 149:6\allowbreak\ 229:7\par
\noindent \textbf{231}: 7:2\allowbreak\ 20:6\allowbreak\ 43:5\allowbreak\ 56:4\allowbreak\ 104:5\allowbreak\ 111:6\allowbreak\ 134:7\allowbreak\ 147:6\allowbreak\ 231:7\par
\noindent \textbf{233}: 7:2\allowbreak\ 9:6\allowbreak\ 18:4\allowbreak\ 43:5\allowbreak\ 45:6\allowbreak\ 70:5\allowbreak\ 81:6\allowbreak\ 106:5\allowbreak\ 109:6\allowbreak\ 134:7\allowbreak\ 136:6\allowbreak\ 170:7\allowbreak\ 172:6\allowbreak\ 197:7\allowbreak\ 208:6\allowbreak\ 233:7\par
\noindent \textbf{234}: 16:6\allowbreak\ 36:5\allowbreak\ 52:6\allowbreak\ 72:5\allowbreak\ 88:6\allowbreak\ 108:5\allowbreak\ 127:7\allowbreak\ 143:6\allowbreak\ 163:7\allowbreak\ 179:6\allowbreak\ 199:7\allowbreak\ 215:6\allowbreak\ 234:7\par
\noindent \textbf{235}: 7:2\allowbreak\ 9:6\allowbreak\ 17:3\allowbreak\ 43:5\allowbreak\ 45:6\allowbreak\ 79:5\allowbreak\ 81:6\allowbreak\ 107:5\allowbreak\ 134:7\allowbreak\ 136:6\allowbreak\ 170:7\allowbreak\ 172:6\allowbreak\ 206:7\allowbreak\ 208:6\allowbreak\ 235:7\par
\noindent \textbf{236}: 7:2\allowbreak\ 18:6\allowbreak\ 43:5\allowbreak\ 45:6\allowbreak\ 54:4\allowbreak\ 79:5\allowbreak\ 81:6\allowbreak\ 106:5\allowbreak\ 134:7\allowbreak\ 145:6\allowbreak\ 170:7\allowbreak\ 172:6\allowbreak\ 206:7\allowbreak\ 208:6\allowbreak\ 236:7\par
\noindent \textbf{238}: 7:2\allowbreak\ 20:6\allowbreak\ 43:5\allowbreak\ 56:4\allowbreak\ 104:5\allowbreak\ 134:7\allowbreak\ 147:6\allowbreak\ 238:7\par
\noindent \textbf{240}: 2:4\allowbreak\ 5:2\allowbreak\ 22:6\allowbreak\ 41:5\allowbreak\ 58:4\allowbreak\ 61:5\allowbreak\ 63:4\allowbreak\ 102:5\allowbreak\ 132:7\allowbreak\ 149:6\allowbreak\ 240:7\par
\noindent \textbf{241}: 2:4\allowbreak\ 3:2\allowbreak\ 4:3\allowbreak\ 5:2\allowbreak\ 22:6\allowbreak\ 23:3\allowbreak\ 41:5\allowbreak\ 58:4\allowbreak\ 61:5\allowbreak\ 63:4\allowbreak\ 101:5\allowbreak\ 132:7\allowbreak\ 149:6\allowbreak\ 241:7\par
\noindent \textbf{243}: 25:4\allowbreak\ 61:5\allowbreak\ 63:4\allowbreak\ 86:6\allowbreak\ 99:5\allowbreak\ 190:7\allowbreak\ 213:6\allowbreak\ 243:7\par
\noindent \textbf{245}: 27:4\allowbreak\ 61:5\allowbreak\ 88:6\allowbreak\ 97:5\allowbreak\ 188:7\allowbreak\ 215:6\allowbreak\ 245:7\par
\noindent \textbf{246}: 2:6\allowbreak\ 24:3\allowbreak\ 28:6\allowbreak\ 36:5\allowbreak\ 38:4\allowbreak\ 60:5\allowbreak\ 61:6\allowbreak\ 63:4\allowbreak\ 64:6\allowbreak\ 96:5\allowbreak\ 127:7\allowbreak\ 129:6\allowbreak\ 151:7\allowbreak\ 155:6\allowbreak\ 187:7\allowbreak\ 188:6\allowbreak\ 190:7\allowbreak\ 191:6\allowbreak\ 246:7\par
\noindent \textbf{247}: 2:6\allowbreak\ 3:2\allowbreak\ 4:3\allowbreak\ 5:2\allowbreak\ 22:6\allowbreak\ 23:3\allowbreak\ 29:6\allowbreak\ 36:5\allowbreak\ 38:4\allowbreak\ 41:5\allowbreak\ 58:4\allowbreak\ 59:5\allowbreak\ 61:6\allowbreak\ 63:4\allowbreak\ 65:6\allowbreak\ 95:5\allowbreak\ 127:7\allowbreak\ 129:6\allowbreak\ 132:7\allowbreak\ 149:6\allowbreak\ 150:7\allowbreak\ 156:6\allowbreak\ 186:7\allowbreak\ 188:6\allowbreak\ 190:7\allowbreak\ 192:6\allowbreak\ 247:7\par
\noindent \textbf{248}: 2:6\allowbreak\ 5:2\allowbreak\ 30:6\allowbreak\ 36:5\allowbreak\ 38:4\allowbreak\ 41:5\allowbreak\ 58:4\allowbreak\ 61:6\allowbreak\ 63:4\allowbreak\ 66:6\allowbreak\ 94:5\allowbreak\ 127:7\allowbreak\ 129:6\allowbreak\ 132:7\allowbreak\ 157:6\allowbreak\ 185:7\allowbreak\ 188:6\allowbreak\ 190:7\allowbreak\ 193:6\allowbreak\ 248:7\par
\noindent \textbf{249}: 1:3\allowbreak\ 6:2\allowbreak\ 31:6\allowbreak\ 42:5\allowbreak\ 57:4\allowbreak\ 62:6\allowbreak\ 63:4\allowbreak\ 67:6\allowbreak\ 93:5\allowbreak\ 133:7\allowbreak\ 158:6\allowbreak\ 184:7\allowbreak\ 189:6\allowbreak\ 190:7\allowbreak\ 194:6\allowbreak\ 249:7\par
\noindent \textbf{250}: 7:2\allowbreak\ 32:6\allowbreak\ 43:5\allowbreak\ 56:4\allowbreak\ 68:6\allowbreak\ 92:5\allowbreak\ 134:7\allowbreak\ 159:6\allowbreak\ 183:7\allowbreak\ 195:6\allowbreak\ 250:7\par
\noindent \textbf{251}: 7:2\allowbreak\ 19:3\allowbreak\ 33:6\allowbreak\ 55:5\allowbreak\ 69:6\allowbreak\ 91:5\allowbreak\ 146:7\allowbreak\ 160:6\allowbreak\ 182:7\allowbreak\ 196:6\allowbreak\ 251:7\par
\noindent \textbf{252}: 7:2\allowbreak\ 9:6\allowbreak\ 18:4\allowbreak\ 34:6\allowbreak\ 43:5\allowbreak\ 45:6\allowbreak\ 54:5\allowbreak\ 81:6\allowbreak\ 90:5\allowbreak\ 134:7\allowbreak\ 136:6\allowbreak\ 145:7\allowbreak\ 161:6\allowbreak\ 170:7\allowbreak\ 172:6\allowbreak\ 181:7\allowbreak\ 208:6\allowbreak\ 252:7\par
\noindent \textbf{253}: 2:4\allowbreak\ 24:3\allowbreak\ 25:6\allowbreak\ 27:4\allowbreak\ 28:6\allowbreak\ 35:4\allowbreak\ 60:5\allowbreak\ 64:6\allowbreak\ 86:5\allowbreak\ 88:6\allowbreak\ 89:5\allowbreak\ 151:7\allowbreak\ 152:6\allowbreak\ 154:7\allowbreak\ 155:6\allowbreak\ 187:7\allowbreak\ 191:6\allowbreak\ 213:7\allowbreak\ 215:6\allowbreak\ 253:7\par
\noindent \textbf{254}: 88:6\allowbreak\ 127:7\allowbreak\ 215:6\allowbreak\ 254:7\par
\noindent \textbf{339}: 6:3\allowbreak\ 20:6\allowbreak\ 23:4\allowbreak\ 26:3\allowbreak\ 172:6\allowbreak\ 302:8\allowbreak\ 316:7\allowbreak\ 322:8\allowbreak\ 339:7\par
\noindent \textbf{341}: 1:3\allowbreak\ 21:7\allowbreak\ 62:6\allowbreak\ 63:4\allowbreak\ 189:6\allowbreak\ 190:7\allowbreak\ 317:8\allowbreak\ 341:7\par
\noindent \textbf{343}: 16:5\allowbreak\ 30:4\allowbreak\ 46:6\allowbreak\ 47:4\allowbreak\ 91:6\allowbreak\ 107:5\allowbreak\ 168:6\allowbreak\ 326:8\allowbreak\ 342:7\allowbreak\ 343:8\par
\noindent \textbf{344}: 40:7\allowbreak\ 42:6\allowbreak\ 82:5\allowbreak\ 133:6\allowbreak\ 167:7\allowbreak\ 169:8\allowbreak\ 173:7\allowbreak\ 175:6\allowbreak\ 209:7\allowbreak\ 302:8\allowbreak\ 342:7\allowbreak\ 344:8\par
\noindent \textbf{345}: 7:2\allowbreak\ 166:7\allowbreak\ 176:8\allowbreak\ 335:7\allowbreak\ 345:8\par
\noindent \textbf{346}: 8:5\allowbreak\ 49:7\allowbreak\ 63:4\allowbreak\ 74:6\allowbreak\ 99:5\allowbreak\ 165:6\allowbreak\ 177:8\allowbreak\ 190:7\allowbreak\ 201:6\allowbreak\ 218:7\allowbreak\ 346:8\par
\noindent \textbf{347}: 9:3\allowbreak\ 37:4\allowbreak\ 46:6\allowbreak\ 51:4\allowbreak\ 161:6\allowbreak\ 164:7\allowbreak\ 330:8\allowbreak\ 342:7\allowbreak\ 347:8\par
\noindent \textbf{348}: 10:6\allowbreak\ 36:5\allowbreak\ 46:6\allowbreak\ 52:5\allowbreak\ 101:6\allowbreak\ 127:7\allowbreak\ 137:6\allowbreak\ 163:7\allowbreak\ 270:8\allowbreak\ 342:7\allowbreak\ 348:8\par
\noindent \textbf{349}: 46:7\allowbreak\ 53:4\allowbreak\ 71:6\allowbreak\ 78:5\allowbreak\ 162:6\allowbreak\ 169:8\allowbreak\ 198:6\allowbreak\ 205:7\allowbreak\ 215:6\allowbreak\ 349:8\par
\noindent \textbf{350}: 9:7\allowbreak\ 17:4\allowbreak\ 26:3\allowbreak\ 54:4\allowbreak\ 161:6\allowbreak\ 169:8\allowbreak\ 178:6\allowbreak\ 350:8\par
\noindent \textbf{351}: 7:2\allowbreak\ 19:3\allowbreak\ 55:7\allowbreak\ 69:6\allowbreak\ 97:5\allowbreak\ 160:6\allowbreak\ 188:8\allowbreak\ 196:6\allowbreak\ 224:7\allowbreak\ 351:8\par
\noindent \textbf{353}: 6:3\allowbreak\ 7:2\allowbreak\ 20:6\allowbreak\ 25:3\allowbreak\ 43:5\allowbreak\ 57:4\allowbreak\ 116:6\allowbreak\ 134:7\allowbreak\ 158:6\allowbreak\ 285:8\allowbreak\ 316:7\allowbreak\ 353:8\par
\noindent \textbf{354}: 10:6\allowbreak\ 16:3\allowbreak\ 36:5\allowbreak\ 46:7\allowbreak\ 58:5\allowbreak\ 107:6\allowbreak\ 127:7\allowbreak\ 137:6\allowbreak\ 157:7\allowbreak\ 205:8\allowbreak\ 215:6\allowbreak\ 276:8\allowbreak\ 326:7\allowbreak\ 354:8\par
\noindent \textbf{357}: 7:2\allowbreak\ 19:3\allowbreak\ 154:7\allowbreak\ 188:8\allowbreak\ 323:7\allowbreak\ 357:8\par
\noindent \textbf{358}: 6:3\allowbreak\ 7:2\allowbreak\ 43:7\allowbreak\ 62:4\allowbreak\ 85:5\allowbreak\ 153:6\allowbreak\ 176:8\allowbreak\ 189:6\allowbreak\ 212:7\allowbreak\ 358:8\par
\noindent \textbf{363}: 4:7\allowbreak\ 6:6\allowbreak\ 21:4\allowbreak\ 42:6\allowbreak\ 67:5\allowbreak\ 133:6\allowbreak\ 148:7\allowbreak\ 169:8\allowbreak\ 173:7\allowbreak\ 300:8\allowbreak\ 342:7\allowbreak\ 363:8\par

\end{multicols}

\end{document}
